\documentclass[10pt]{article}
\usepackage[margin=0.85in]{geometry}
\usepackage[T1]{fontenc}
\usepackage{lmodern}
\usepackage{amsmath,amssymb,amsthm}
\usepackage{booktabs}
\usepackage{microtype}
\usepackage[hidelinks]{hyperref}
\newtheorem{theorem}{Theorem}
\newtheorem{proposition}{Proposition}
\newcommand{\Z}{\mathbb Z}
\newcommand{\Hom}{\operatorname{Hom}}
\newcommand{\diag}{\operatorname{diag}}
\title{Splitting a Concrete Defect Extension in a Restricted Mark Cokernel}
\author{Anonymous}
\date{}
\begin{document}
\maketitle

\begin{abstract}
For a frozen $16\times16$ restricted self-mark matrix $M$, earlier exact
calculations locate a subgroup
$D\cong\mathbb Z/8\oplus(\mathbb Z/2)^2$ inside
$C=\mathbb Z^{16}/M\mathbb Z^{16}$.  Abstract invariant factors do not decide
whether this particular inclusion splits.  We give a three-coordinate
retraction, an explicit index-$32$ lattice for its kernel, and an integral
presentation of the resulting complement.  We also prove that the fixed
subgroup has exactly $2^{41}$ complements.  All claims concern only the named
sixteen-type support.
\end{abstract}

\section{The concrete extension}
Let $M$ be the exact restricted mark matrix, let
$U=[u_1\ u_2\ u_3]$, and put
\[
 C=\Z^{16}/M\Z^{16},\qquad
 A=\Z/8\oplus\Z/2\oplus\Z/2.
\]
The preceding certificate proves that $a\mapsto[Ua]$ embeds $A$ in $C$;
write its image as $D$.  It also gives
\[
 C/D\cong(\Z/2)^8\oplus(\Z/4)^2\oplus\Z/12\oplus\Z/144.
\]
The question here is about the actual extension
$0\to D\to C\to C/D\to0$.  The abstract groups alone are insufficient:
$0\to\Z/2\to\Z/4\to\Z/2\to0$, for example, does not split.

The scope firewall is
\[
\texttt{NO\_BAD\_EULER\_OR\_ROOT\_NUMBER}.
\]
No full table of marks, full Burnside ring, arithmetic or local construction,
Euler factor, root number, automorphy, or Hilbert--Polya operator is claimed.

\section{A sparse retraction}
For $x=(x_1,\ldots,x_{16})^T$, define
\[
 \rho([x])=(x_{10}\bmod8,\ x_3\bmod2,\ x_1+x_{15}\bmod2).
\]
Its integer row matrix is
\[
 R=\begin{pmatrix}
 e_{10}^T\\ e_3^T\\ e_1^T+e_{15}^T
 \end{pmatrix},\qquad \Delta=\diag(8,2,2).
\]

\begin{theorem}[Splitting of the actual inclusion]
The map $\rho:C\to A$ is well defined and its restriction to $D$ is the
identity.  Consequently $C=D\oplus\ker\rho$.
\end{theorem}

\begin{proof}
Exact integer multiplication gives $RM\equiv0$ rowwise modulo $(8,2,2)$,
so $\rho$ descends through $M\Z^{16}$.  It also gives
\[
 RU=\diag(1,-1,-1)\equiv I_3\pmod\Delta.
\]
Thus $\rho([Ua])=a$ in $A$ for every $a$, which is precisely the retraction
identity for the fixed inclusion.
\end{proof}

\section{An explicit complement}
Let
\[
 L=\{x\in\Z^{16}:x_{10}\equiv0\ (8),\ x_3\equiv0\ (2),\
 x_1+x_{15}\equiv0\ (2)\}.
\]
This is the full inverse image of zero under the congruence map represented by
$R$; in particular, it is not the integer kernel of $R$.  A column basis $B$
is obtained from the standard basis by replacing
\[
 b_1=e_1-e_{15},\quad b_3=2e_3,\quad
 b_{10}=8e_{10},\quad b_{15}=2e_{15}.
\]
It has determinant $32$.  Since $RM\equiv0\pmod\Delta$,
$M\Z^{16}\subset L$.

\begin{proposition}[Complement presentation]
Put $K=L/M\Z^{16}$.  The matrix $N=B^{-1}M$ is integral, satisfies $BN=M$,
and has Smith invariants
\[
 (1,1,1,1,2,2,2,2,2,2,2,2,4,4,12,144).
\]
Hence
\[
 K\cong(\Z/2)^8\oplus(\Z/4)^2\oplus\Z/12\oplus\Z/144,
 \qquad |K|=7077888.
\]
\end{proposition}

\begin{proof}
The displayed basis describes the three congruences directly.  Exact
Gauss--Jordan elimination makes $N$ integral; Euclidean Smith reduction gives
the stated diagonal.  The product check
$|D||K|=32\cdot7077888=226492416=|C|$ is independent confirmation.
\end{proof}

\section{Counting every complement}
Fix the displayed retraction $\rho_0$.  Any other retraction differs from it
by a homomorphism vanishing on $D$, hence by a unique member of
$\Hom(C/D,D)$.  Conversely every such homomorphism, composed with the quotient
map and added to $\rho_0$, is a retraction.  Thus the retractions form a
torsor under $\Hom(C/D,D)$.

The $2$-primary part of the quotient is
\[
 (C/D)_{(2)}\cong(\Z/2)^8\oplus(\Z/4)^3\oplus\Z/16.
\]
Using $|\Hom(\Z/2^a,\Z/2^b)|=2^{\min(a,b)}$, the exponents for the three
target factors $\Z/8,\Z/2,\Z/2$ are respectively $17,12,12$.  Therefore
\[
 |\Hom(C/D,D)|=2^{17+12+12}=2^{41}=2199023255552.
\]
For a fixed inclusion, taking kernels bijects retractions with complements:
a chosen complement determines the unique projection that is the identity on
$D$.  Hence there are exactly $2^{41}$ complements.  In particular, the
explicit $K$ above is not canonical.

\section{Reproducibility and conclusion}
The evidence binds the exact C64, C65, C66, and C68 bytes.  Producer and
checker independently reconstruct $R,B,N$, the Smith diagonal, and the
homomorphism count; a separate SymPy path reproduces the matrix identities and
Smith form.  Clean replay passes and twenty-three hostile semantic mutations
are rejected.  C69 therefore resolves the extension question constructively:
the concrete defect subgroup splits, its displayed complement has the C68
quotient invariants, and all complements are counted exactly.

\end{document}
